\chapter{Arnoldi and Hessenberg Decomposition}\label{v0ch:11}
\begin{chaptersummary}
Krylov powers generate a search space, but increasing linear dependence complicates numerical use. Arnoldi orthogonalizes each new direction and represents the large operator in a small Hessenberg matrix. This verifiable connection between large and small systems supplies the least-squares problem used by GMRES in the next chapter.
\end{chaptersummary}
\learninggoals{Perform MGS Arnoldi; calculate two steps; explain loss of orthogonality and happy breakdown.}
\paragraph{Core Question}How can a reliable subspace basis be constructed without storing high matrix powers?

\section{The Arnoldi Recurrence}
Set $\beta=\norm{\bm r_0}_2$, $\bm v_1=\bm r_0/\beta$. At step $j$,
\[
 \bm w=A\bm v_j;\quad
 h_{ij}=\bm v_i^T\bm w,\quad
 \bm w\leftarrow\bm w-h_{ij}\bm v_i\quad(i=1,\ldots,j);
 \qquad h_{j+1,j}=\norm{\bm w}_2.
\]
If $h_{j+1,j}\ne0$, set $\bm v_{j+1}=\bm w/h_{j+1,j}$. Consequently,
\[
 AV_m=V_{m+1}\bar H_m,\qquad V_{m+1}^TV_{m+1}=I
\]
in exact arithmetic, where $\bar H_m$ is $(m+1)\times m$ and upper Hessenberg.
\section{Two Steps for the Four-by-Four Matrix}
Use $A$ from Chapter 4 and $\bm r_0=\bm e_1$. First,
\[
 A\bm e_1=(0,4,2,0)^T,\quad h_{11}=0,\quad
 h_{21}=2\sqrt5,\quad
 \bm v_2=(0,2,1,0)^T/\sqrt5.
\]
Next, $A\bm v_2=(5,8,3,3)^T/\sqrt5$. Sequential projection gives
\[
 h_{12}=\sqrt5,\quad h_{22}=19/5,\quad
 h_{32}=7/5,\quad
 \bm v_3=(0,2,-4,15)^T/(7\sqrt5),
\]
\[
 \bar H_2=\begin{bmatrix}0&\sqrt5\\2\sqrt5&19/5\\0&7/5\end{bmatrix}.
\]
Check $\bm v_i^T\bm v_j=\delta_{ij}$ and $AV_2=V_3\bar H_2$.
\conceptbridge{How does a small matrix retain solution information from the original problem?}{Arnoldi represents the action on generated basis vectors exactly in the extended basis. The later approximation comes from restricting the trial space, not from approximating these represented actions. Finite precision also requires monitoring orthogonality.}{Large operator action $\longrightarrow$ orthogonal coordinates $\longrightarrow$ small Hessenberg problem.}
\section{Breakdown Diagnostics and Cost}
The condition $h_{j+1,j}=0$ identifies an $A$-invariant subspace. For an invertible system and the current residual, residual minimization then gives the exact solution: happy breakdown. In finite precision, a nearly zero subdiagonal instead triggers a true-residual check against the convergence tolerance.
\costnote{$m$ steps require approximately $m$ matrix products, $O(Nm^2)$ orthogonalization work, and $O(Nm)$ storage. MGS can lose orthogonality; repeat projection or use a more stable strategy when needed. Distributed inner products incur reductions.}
\coderef{sparse_linear_algebra/solvers.py}
\aiexercise{Recompute both steps in MGS order, compare $AV_2$ with $V_3\bar H_2$, and verify every inner product.}{Perturb a power basis by roundoff and compare orthogonality after classical Gram--Schmidt, MGS, and two-pass MGS.}
\studyroute{Selective reorthogonalization and communication-avoiding Arnoldi are omitted. See \cite{saad}.}{Arnoldi; Hessenberg; reorthogonalization; happy breakdown; orthogonality loss}
